\chapter{Formal Specification of the RSVP Field System}
\label{app:rsvp}

\epigraph{The universe behaves less like a completed geometric sculpture and more
like an ongoing constraint propagation process in which prior configurations
continuously reshape future reachable manifolds.}{}

\noindent The Relativistic Scalar-Vector Plenum (\RSVP) is a coupled field system
providing the physical and mathematical substrate on which the Yarncrawler framework
is grounded. This appendix develops the formal specification from first principles,
states the seven governing axioms, presents the Lagrangian structure, and derives
the field equations. Claim-status throughout: the cosmological applications are
\textbf{[C]}; the field formalism as a modeling language for constraint propagation
is \textbf{[H]}.

\section{Field Definitions}

\begin{definition}{RSVP Field Triple}
The \RSVP field system is a triple $(\Phi, \mathbf{v}, S)$ of dynamical fields
over a spacetime manifold $\mathcal{M}$:
\begin{itemize}
  \item $\Phi(\mathbf{x},t) \in \mathbb{R}$: the \emph{scalar density field},
    encoding concentration, potential, legitimacy, or local capacity at each
    spacetime point. High $\Phi$ corresponds to dense, stable, or energetically
    accessible regions.
  \item $\mathbf{v}(\mathbf{x},t) \in T_\mathbf{x}\mathcal{M}$: the \emph{vector
    flow field}, encoding directed transport, trajectories, and causal propagation.
    Flows along $\mathbf{v}$ define the admissible directions of system evolution.
  \item $S(\mathbf{x},t) \in \mathbb{R}_{\geq 0}$: the \emph{entropy field},
    encoding disorder, unresolved configurational accessibility, and complexity
    budget. $S$ functions as the reservoir of unresolved structure available for
    future reinterpretation.
\end{itemize}
\end{definition}

\section{The Seven Axioms}

\begin{axiom}{Non-negativity and Boundedness of Entropy}
$S(\mathbf{x},t) \geq 0$ everywhere and for all $t$. Entropy cannot be negative.
In regions of complete constraint satisfaction, $S \to 0$; in regions of maximal
configurational uncertainty, $S$ is bounded by the local informational capacity
of the field.
\end{axiom}

\begin{axiom}{Flow-Density Coupling}
The vector field $\mathbf{v}$ is coupled to the scalar field $\Phi$ through
the continuity relation:
\[
  \frac{\partial \Phi}{\partial t} + \nabla \cdot (\Phi\, \mathbf{v}) = \sigma_\Phi,
\]
where $\sigma_\Phi$ is a source/sink term encoding the creation or destruction of
scalar density through constraint events. In the absence of constraint events,
$\sigma_\Phi = 0$ and scalar density is conserved along flow lines.
\end{axiom}

\begin{axiom}{Entropy-Driven Flow}
The vector field evolves under the influence of entropy gradients:
\[
  \frac{\partial \mathbf{v}}{\partial t} + (\mathbf{v}\cdot\nabla)\mathbf{v}
  = -\nabla S + \mathbf{f}_{\rm ext},
\]
where $\mathbf{f}_{\rm ext}$ encodes external constraint forces. Flows are
driven from high-entropy regions toward low-entropy regions: the system
navigates toward configurations of greater constraint stability.
\end{axiom}

\begin{axiom}{Entropy Production and Relaxation}
Entropy evolves through production and relaxation:
\[
  \frac{\partial S}{\partial t} + \nabla \cdot (S\,\mathbf{v})
  = \Pi_S - \Lambda_S,
\]
where $\Pi_S \geq 0$ is the entropy production rate (irreversibility) and
$\Lambda_S \geq 0$ is the entropy relaxation rate (constraint stabilization).
The second law requires $\Pi_S \geq \Lambda_S$ globally, but local relaxation
is possible at the cost of increased entropy elsewhere.
\end{axiom}

\begin{axiom}{Admissibility Constraint on Trajectories}
A trajectory $\gamma(t)$ through the \RSVP field manifold is admissible if and
only if it satisfies the local compatibility conditions imposed by all three
fields simultaneously:
\[
  \mathbf{v}(\gamma(t),t) \cdot \nabla S(\gamma(t),t) \leq 0,
\]
that is, admissible trajectories do not move against the entropy gradient. They
flow from higher to lower local entropy, or along isoentropic surfaces.
\end{axiom}

\begin{axiom}{Structural Residue Accumulation}
The fields accumulate structural residue across time: prior field configurations
influence future evolution through path-dependent coupling terms. Define the
residue field:
\[
  \mathcal{R}(\mathbf{x},t) = \int_{-\infty}^{t} K(t-t')\,
  \Phi(\mathbf{x},t')\,dt',
\]
where $K(t-t')$ is a memory kernel with $K(\tau) \to 0$ as $\tau \to \infty$.
The residue field enters the Lagrangian as a history-dependent potential,
ensuring that the present state inherits structure from prior configurations.
\end{axiom}

\begin{axiom}{Semantic Interpretability}
Under the semantic interpretation developed in Chapter~\ref{ch:yarncrawler_formal},
the \RSVP fields map to:
\begin{itemize}
  \item $\Phi$: semantic density, knowledge concentration, or accumulated
    legitimacy in a representational system.
  \item $\mathbf{v}$: recursive trajectories of meaning, directed inference
    flows, or causal propagation of conceptual updates.
  \item $S$: unresolved noise available for future reinterpretation; the
    non-trivial cohomological residue of the sheaf-theoretic framework.
\end{itemize}
This semantic interpretation does not require the cosmological interpretation
to be valid. The field formalism functions as a modeling language independently
of its physical instantiation claims.
\end{axiom}

\section{The Lagrangian Structure}

The \RSVP dynamics are generated by the Lagrangian density:
\[
  \mathcal{L}_{\rm RSVP} = \mathcal{L}_\Phi + \mathcal{L}_v + \mathcal{L}_S
  + \mathcal{L}_{\rm int},
\]
where the component terms are:
\begin{align}
  \mathcal{L}_\Phi &= \frac{1}{2}(\partial_t \Phi)^2 - \frac{c_\Phi^2}{2}|\nabla\Phi|^2
    - V(\Phi), \\
  \mathcal{L}_v &= \frac{\rho_v}{2}|\partial_t \mathbf{v}|^2
    - \frac{\mu_v}{2}|\nabla \times \mathbf{v}|^2, \\
  \mathcal{L}_S &= -S\log S + \frac{\kappa_S}{2}|\nabla S|^2, \\
  \mathcal{L}_{\rm int} &= -\lambda_{\Phi S}\,\Phi\,S
    - \lambda_{vS}\,\mathbf{v}\cdot\nabla S
    - \lambda_{\mathcal{R}}\,\mathcal{R}\,\Phi.
\end{align}
The coupling constants $\lambda_{\Phi S}$, $\lambda_{vS}$, $\lambda_{\mathcal{R}}$
govern the interaction strengths between the three fields and the residue term.
The potential $V(\Phi)$ may take various forms depending on the application
domain; for the cosmological interpretation, $V(\Phi) = \frac{m^2}{2}\Phi^2 +
\frac{\lambda}{4}\Phi^4$ is a symmetry-breaking potential.

Variation of $\mathcal{L}_{\rm RSVP}$ with respect to each field yields the
Euler-Lagrange equations, which reduce to the axiom-specified evolution equations
under the appropriate coupling limits.

\section{Cosmological Application}
\label{app:rsvp:cosmo}

\claimC Under the cosmological interpretation, expansion is replaced by entropic
relaxation. The scalar field $\Phi$ encodes matter density; the entropy field $S$
replaces the scale factor as the primary measure of cosmic evolution; and $\mathbf{v}$
encodes the peculiar velocity field of matter.

Cosmological redshift on this interpretation arises from photons traversing
entropy gradients:
\[
  z_{\rm RSVP} \approx \int_0^L \kappa\,|\nabla S(\ell)|\,d\ell,
\]
where $\kappa$ is a coupling constant relating entropy gradient to energy loss
and $L$ is the path length. This formula makes distinct observational predictions
from $\Lambda$CDM at $k \gtrsim 0.1\,h/\text{Mpc}$ and at the Silk damping scale.

\claimC Gravity on this interpretation emerges as entropy descent: the
gravitational acceleration of a test body is proportional to $-\nabla S$ at
leading order, so that matter falls toward entropy minima rather than toward
mass concentrations in the conventional sense. The connection to the MOND
phenomenological transition---the observed departure from Newtonian dynamics
at low accelerations---is that the entropy gradient structure of the $S$ field
exhibits a characteristic scale at which the gradient profile transitions from
the high-acceleration regime (tracking baryonic density closely) to the
low-acceleration regime (governed by the large-scale entropy field topology).

All cosmological claims in this appendix carry status \textbf{[C]} and require
independent observational validation. The field formalism itself is independent
of these applications.

\section{Derived Stacks and the Four Projection Modes}

The companion paper \emph{Axioms for a Falling Universe} develops four derived
stacks from the \RSVP Lagrangian, corresponding to four projection modes of the
field system onto different observable domains:
\begin{enumerate}
  \item \textbf{Thermodynamic stack}: projection onto entropy production and
    relaxation dynamics; yields predictions about large-scale structure formation
    timescales.
  \item \textbf{Kinetic stack}: projection onto flow field statistics; yields
    predictions about peculiar velocity distributions.
  \item \textbf{Informational stack}: projection onto the residue field; yields
    predictions about persistence timescales and structural memory.
  \item \textbf{Semantic stack}: projection onto the symbolic interpretation;
    yields the Yarncrawler dynamics of Chapter~\ref{ch:yarncrawler_formal}.
\end{enumerate}
The four stacks are not independent theories but coordinate representations of
the same underlying field system, related by the transition maps of the semantic
atlas introduced in Chapter~\ref{ch:manifold}.
